
\documentclass{article}
\usepackage{colortbl}
\usepackage{makecell}
\usepackage{multirow}
\usepackage{supertabular}

\begin{document}

\newcounter{utterance}

\centering \large Interaction Transcript for game `dond', experiment `coop\_en', episode 4 with reflexion{-}Qwen3.5{-}27B.
\vspace{24pt}

{ \footnotesize  \setcounter{utterance}{1}
\setlength{\tabcolsep}{0pt}
\begin{supertabular}{c@{$\;$}|p{.15\linewidth}@{}p{.15\linewidth}p{.15\linewidth}p{.15\linewidth}p{.15\linewidth}p{.15\linewidth}}
    \# & $\;$A & \multicolumn{4}{c}{Game Master} & $\;\:$B\\
    \hline

    \theutterance \stepcounter{utterance}  
    & & \multicolumn{4}{p{0.6\linewidth}}{
        \cellcolor[rgb]{0.9,0.9,0.9}{
            \makecell[{{p{\linewidth}}}]{
                \texttt{\tiny{[P1$\langle$GM]}}
                \texttt{You are playing a cooperative negotiation game in which you have to agree on how to divide a set of items among you and another player.} \\
\\ 
\texttt{Rules:} \\
\texttt{(a) You and the other player are given a set of items. Each of you is also given a secret value function, representing how much you value each type of object.} \\
\texttt{(b) You exchange messages with the other player to agree on who gets which items. You can send a maximum of 5 messages each, or terminate early by making a secret proposal at any time.} \\
\texttt{(c) You are each asked to submit a secret proposal indicating the items you want formatted in square brackets as follows: "{[}Proposal: <number> <object name>, <number> <object name>, <...>{]}"} \\
\texttt{(d) If your proposals are complementary, i.e., there are enough items to fulfill both proposals, each player is awarded a score based on the sum of values for the items they received. Otherwise, both of you get zero points.} \\
\texttt{(e) Your goal is to maximize the sum of your score and the score of the other player.} \\
\\ 
\texttt{Let us start.} \\
\\ 
\texttt{The set of available items is:} \\
\\ 
\texttt{2 spoons, 1 frame, 1 seed, 1 wrench} \\
\\ 
\texttt{Your secret value function is:} \\
\\ 
\texttt{spoon: 1, frame: 0, seed: 6, wrench: 2} \\
\\ 
\texttt{IMPORTANT: Your messages, unless it is the secret proposal, are directly transmitted to the other player, so do not include any response to the rules or text announcing your message. To make a secret proposal, use the indicated format. Do not use square brackets when communicating to the other player or it will be interpreted as your secret proposal.} \\
            }
        }
    }
    & \\ \\

    \theutterance \stepcounter{utterance}  
    & \multicolumn{4}{p{0.6\linewidth}}{
        \cellcolor[rgb]{0.9,0.9,0.9}{
            \makecell[{{p{\linewidth}}}]{
                \texttt{\tiny{[P1$\rangle$GM]}}
                \texttt{Hi! Let's figure out the best way to split these items. I value the seed the most (6 points), followed by the wrench (2 points), and the spoons are worth 1 each. I don't really care about the frame (0 points).} \\
\\ 
\texttt{What are your values for these items? Knowing that will help us make sure we both get the things we want most.} \\
            }
        }
    }
    & & \\ \\

    \theutterance \stepcounter{utterance}  
    & & & \multicolumn{2}{p{0.3\linewidth}}{
        \cellcolor[rgb]{0.9,0.9,0.9}{
            \makecell[{{p{\linewidth}}}]{
                \texttt{\tiny{[GM$|$GM]}}
                \texttt{continue} \\
            }
        }
    }
    & & \\ \\

    \theutterance \stepcounter{utterance}  
    & & & \multicolumn{2}{p{0.3\linewidth}}{
        \cellcolor[rgb]{0.9,0.9,0.9}{
            \makecell[{{p{\linewidth}}}]{
                \texttt{\tiny{[GM$|$GM]}}
                \texttt{Hi! Let's figure out the best way to split these items. I value the seed the most (6 points), followed by the wrench (2 points), and the spoons are worth 1 each. I don't really care about the frame (0 points).} \\
\\ 
\texttt{What are your values for these items? Knowing that will help us make sure we both get the things we want most.} \\
            }
        }
    }
    & & \\ \\

    \theutterance \stepcounter{utterance}  
    & & \multicolumn{4}{p{0.6\linewidth}}{
        \cellcolor[rgb]{0.9,0.9,0.9}{
            \makecell[{{p{\linewidth}}}]{
                \texttt{\tiny{[GM$\rangle$P2]}}
                \texttt{You are playing a cooperative negotiation game in which you have to agree on how to divide a set of items among you and another player.} \\
\\ 
\texttt{Rules:} \\
\texttt{(a) You and the other player are given a set of items. Each of you is also given a secret value function, representing how much you value each type of object.} \\
\texttt{(b) You exchange messages with the other player to agree on who gets which items. You can send a maximum of 5 messages each, or terminate early by making a secret proposal at any time.} \\
\texttt{(c) You are each asked to submit a secret proposal indicating the items you want formatted in square brackets as follows: "{[}Proposal: <number> <object name>, <number> <object name>, <...>{]}"} \\
\texttt{(d) If your proposals are complementary, i.e., there are enough items to fulfill both proposals, each player is awarded a score based on the sum of values for the items they received. Otherwise, both of you get zero points.} \\
\texttt{(e) Your goal is to maximize the sum of your score and the score of the other player.} \\
\\ 
\texttt{Let us start.} \\
\\ 
\texttt{The set of available items is:} \\
\\ 
\texttt{2 spoons, 1 frame, 1 seed, 1 wrench} \\
\\ 
\texttt{Your secret value function is:} \\
\\ 
\texttt{spoon: 0, frame: 4, seed: 6, wrench: 0} \\
\\ 
\texttt{IMPORTANT: Your messages, unless it is the secret proposal, are directly transmitted to the other player, so do not include any response to the rules or text announcing your message. To make a secret proposal, use the indicated format. Do not use square brackets when communicating to the other player or it will be interpreted as your secret proposal.} \\
\\ 
\\ 
\texttt{Hi! Let's figure out the best way to split these items. I value the seed the most (6 points), followed by the wrench (2 points), and the spoons are worth 1 each. I don't really care about the frame (0 points).} \\
\\ 
\texttt{What are your values for these items? Knowing that will help us make sure we both get the things we want most.} \\
            }
        }
    }
    & \\ \\

    \theutterance \stepcounter{utterance}  
    & & & \multicolumn{4}{p{0.6\linewidth}}{
        \cellcolor[rgb]{0.9,0.9,0.9}{
            \makecell[{{p{\linewidth}}}]{
                \texttt{\tiny{[GM$\langle$P2]}}
                \texttt{Hi! Thanks for sharing your values. To help us maximize our combined score, here are mine:} \\
\\ 
\texttt{I value the **seed** the most (6 points), followed by the **frame** (4 points). The **spoons** and the **wrench** are worth 0 to me.} \\
\\ 
\texttt{Since we both value the seed highly, we need to decide who gets it. However, I see a clear opportunity for efficiency:} \\
\texttt{{-} You value the wrench (2 points) and spoons (1 each), while I value them at 0.} \\
\texttt{{-} I value the frame (4 points), while you value it at 0.} \\
\texttt{{-} We both value the seed (6 points).} \\
\\ 
\texttt{If I take the **frame** and you take the **wrench** and **spoons**, we both get items we value. The only conflict is the **seed**. Since we both value it equally at 6, giving it to either of us adds 6 to the total score.} \\
\\ 
\texttt{How about this plan:} \\
\texttt{{-} **You get**: 1 seed, 1 wrench, 2 spoons.} \\
\texttt{{-} **I get**: 1 frame.} \\
\\ 
\texttt{This way, you get your top priority (seed) plus the wrench and spoons you like, and I get the frame which is my second highest priority. This seems to maximize our total points without conflict. Does this work for you?} \\
            }
        }
    }
     \\ \\

    \theutterance \stepcounter{utterance}  
    & & & \multicolumn{2}{p{0.3\linewidth}}{
        \cellcolor[rgb]{0.9,0.9,0.9}{
            \makecell[{{p{\linewidth}}}]{
                \texttt{\tiny{[GM$|$GM]}}
                \texttt{continue} \\
            }
        }
    }
    & & \\ \\

    \theutterance \stepcounter{utterance}  
    & & & \multicolumn{2}{p{0.3\linewidth}}{
        \cellcolor[rgb]{0.9,0.9,0.9}{
            \makecell[{{p{\linewidth}}}]{
                \texttt{\tiny{[GM$|$GM]}}
                \texttt{Hi! Thanks for sharing your values. To help us maximize our combined score, here are mine:} \\
\\ 
\texttt{I value the **seed** the most (6 points), followed by the **frame** (4 points). The **spoons** and the **wrench** are worth 0 to me.} \\
\\ 
\texttt{Since we both value the seed highly, we need to decide who gets it. However, I see a clear opportunity for efficiency:} \\
\texttt{{-} You value the wrench (2 points) and spoons (1 each), while I value them at 0.} \\
\texttt{{-} I value the frame (4 points), while you value it at 0.} \\
\texttt{{-} We both value the seed (6 points).} \\
\\ 
\texttt{If I take the **frame** and you take the **wrench** and **spoons**, we both get items we value. The only conflict is the **seed**. Since we both value it equally at 6, giving it to either of us adds 6 to the total score.} \\
\\ 
\texttt{How about this plan:} \\
\texttt{{-} **You get**: 1 seed, 1 wrench, 2 spoons.} \\
\texttt{{-} **I get**: 1 frame.} \\
\\ 
\texttt{This way, you get your top priority (seed) plus the wrench and spoons you like, and I get the frame which is my second highest priority. This seems to maximize our total points without conflict. Does this work for you?} \\
            }
        }
    }
    & & \\ \\

    \theutterance \stepcounter{utterance}  
    & & \multicolumn{4}{p{0.6\linewidth}}{
        \cellcolor[rgb]{0.9,0.9,0.9}{
            \makecell[{{p{\linewidth}}}]{
                \texttt{\tiny{[P1$\langle$GM]}}
                \texttt{Hi! Thanks for sharing your values. To help us maximize our combined score, here are mine:} \\
\\ 
\texttt{I value the **seed** the most (6 points), followed by the **frame** (4 points). The **spoons** and the **wrench** are worth 0 to me.} \\
\\ 
\texttt{Since we both value the seed highly, we need to decide who gets it. However, I see a clear opportunity for efficiency:} \\
\texttt{{-} You value the wrench (2 points) and spoons (1 each), while I value them at 0.} \\
\texttt{{-} I value the frame (4 points), while you value it at 0.} \\
\texttt{{-} We both value the seed (6 points).} \\
\\ 
\texttt{If I take the **frame** and you take the **wrench** and **spoons**, we both get items we value. The only conflict is the **seed**. Since we both value it equally at 6, giving it to either of us adds 6 to the total score.} \\
\\ 
\texttt{How about this plan:} \\
\texttt{{-} **You get**: 1 seed, 1 wrench, 2 spoons.} \\
\texttt{{-} **I get**: 1 frame.} \\
\\ 
\texttt{This way, you get your top priority (seed) plus the wrench and spoons you like, and I get the frame which is my second highest priority. This seems to maximize our total points without conflict. Does this work for you?} \\
            }
        }
    }
    & \\ \\

    \theutterance \stepcounter{utterance}  
    & \multicolumn{4}{p{0.6\linewidth}}{
        \cellcolor[rgb]{0.9,0.9,0.9}{
            \makecell[{{p{\linewidth}}}]{
                \texttt{\tiny{[P1$\rangle$GM]}}
                \texttt{That sounds like a perfect plan! Since we both value the seed equally, giving it to you or me doesn't change the total combined score, but giving it to me along with the wrench and spoons ensures I get my high{-}value items, while you get the frame which is your top priority.} \\
\\ 
\texttt{Let's confirm the split:} \\
\texttt{{-} **You**: 1 frame} \\
\texttt{{-} **Me**: 1 seed, 1 wrench, 2 spoons} \\
\\ 
\texttt{This leaves no items unassigned and maximizes our joint score. I agree to this proposal.} \\
\\ 
\texttt{{[}Proposal: 1 seed, 1 wrench, 2 spoons{]}} \\
            }
        }
    }
    & & \\ \\

    \theutterance \stepcounter{utterance}  
    & & & \multicolumn{2}{p{0.3\linewidth}}{
        \cellcolor[rgb]{0.9,0.9,0.9}{
            \makecell[{{p{\linewidth}}}]{
                \texttt{\tiny{[GM$|$GM]}}
                \texttt{proposal} \\
            }
        }
    }
    & & \\ \\

    \theutterance \stepcounter{utterance}  
    & & & \multicolumn{2}{p{0.3\linewidth}}{
        \cellcolor[rgb]{0.9,0.9,0.9}{
            \makecell[{{p{\linewidth}}}]{
                \texttt{\tiny{[GM$|$GM]}}
                [2, 0, 1, 1]
            }
        }
    }
    & & \\ \\

    \theutterance \stepcounter{utterance}  
    & & \multicolumn{4}{p{0.6\linewidth}}{
        \cellcolor[rgb]{0.9,0.9,0.9}{
            \makecell[{{p{\linewidth}}}]{
                \texttt{\tiny{[GM$\rangle$P2]}}
                \texttt{The other player has submitted a secret proposal. It is now your turn to submit a proposal of your own.} \\
            }
        }
    }
    & \\ \\

    \theutterance \stepcounter{utterance}  
    & & & \multicolumn{4}{p{0.6\linewidth}}{
        \cellcolor[rgb]{0.9,0.9,0.9}{
            \makecell[{{p{\linewidth}}}]{
                \texttt{\tiny{[GM$\langle$P2]}}
                \texttt{{[}Proposal: 1 frame{]}} \\
            }
        }
    }
     \\ \\

    \theutterance \stepcounter{utterance}  
    & & & \multicolumn{2}{p{0.3\linewidth}}{
        \cellcolor[rgb]{0.9,0.9,0.9}{
            \makecell[{{p{\linewidth}}}]{
                \texttt{\tiny{[GM$|$GM]}}
                \texttt{proposal} \\
            }
        }
    }
    & & \\ \\

    \theutterance \stepcounter{utterance}  
    & & & \multicolumn{2}{p{0.3\linewidth}}{
        \cellcolor[rgb]{0.9,0.9,0.9}{
            \makecell[{{p{\linewidth}}}]{
                \texttt{\tiny{[GM$|$GM]}}
                [0, 1, 0, 0]
            }
        }
    }
    & & \\ \\

    \theutterance \stepcounter{utterance}  
    & & & \multicolumn{2}{p{0.3\linewidth}}{
        \cellcolor[rgb]{0.9,0.9,0.9}{
            \makecell[{{p{\linewidth}}}]{
                \texttt{\tiny{[GM$|$GM]}}
                [[2, 0, 1, 1], [0, 1, 0, 0]]
            }
        }
    }
    & & \\ \\

\end{supertabular}
}

\end{document}
